\section{Evaluation Questions and Evidence Boundaries}\label{sec:evaluation}
The study addresses four questions with different units of evidence.
First, which cooperative phase profiles resist unilateral deviations in
the specified game? Second, how do entangler families change payoffs and
fairness at a fixed strategy? Third, which apparent robustness differences
survive controls for circuit implementation? Fourth, which ideal properties
remain supportable on a finite noisy processor? Table~\ref{tab:evidence-boundaries}
states the operational test and interpretation of each question.

\begin{table*}[t]
\centering
\caption{Evaluation questions, evidence layers, and interpretation boundaries.}
\label{tab:evidence-boundaries}
\begin{tabular}{p{.19\textwidth}p{.25\textwidth}p{.23\textwidth}p{.23\textwidth}}
\toprule
Question & Operational test & Evidence & Interpretation permitted\\
\midrule
Which phases support cooperation? & Compare every named unilateral strategy
with the cooperative payoff. & Phase-boundary proof, unrestricted witness,
and dense-circuit cross-checks. & Exact ideal statements for the specified
operator and menu; no unrestricted equilibrium claim.\\
How does the entangler affect payoffs? & Evaluate the same fixed strategy
and enumerate finite-menu best replies. & Five-family noiseless landscape,
per-player vectors, and graph-symmetry checks. & A fixed-strategy comparison;
not each family's globally optimal play.\\
What determines noise robustness? & Compare payoff gaps and equilibrium
status; vary compilation and match gate counts. & Gate-level sweeps,
implementation controls, and wiring permutations. & Production-circuit
robustness; geometry and synthesis are not interchangeable.\\
What survives on hardware? & Test all-player deviation gaps under a frozen
confidence rule; inspect welfare and population separately. & Registered
circuits, raw counts, calibration records, and repeated batches. & Bounded
implemented-game evidence; no distributed-market deployment claim.\\
\bottomrule
\end{tabular}
\end{table*}

\subsection{Mathematical identities and empirical estimands}
The phase boundary and four-product-state representation are algebraic
properties of the ideal GHZ protocol. Finite floating-point tests check
their implementation but do not establish the all-$N$ statements by
themselves. Hardware probabilities instead estimate repeated circuit
outcomes. The mean payoff, each player's payoff, the worst unilateral
gain, and the all-zero population are distinct estimands. In particular,
high mean payoff cannot substitute for a player-level incentive test.

\subsection{Evidence hierarchy and negative findings}
Historical runs, retrospective reweighting, and prospectively registered
tests retain their separate status. Failed noise predictions remain in
the record. The alternative-phase repeat is judged under the same rule
as the first batch, without pooling to reverse a failed acceptance test.
No result on a finite device is promoted to a statement about unlimited
physical scaling. Compact classical evaluation is reported as an analysis
tool, while physical player count remains the number of simultaneously
represented qubits in the executed protocol.
